% !TeX root = ../../poisson-cubic.tex
% Bibliography for the Poisson manuscript.
\begin{thebibliography}{99}
\addcontentsline{toc}{section}{参考文献}
\small
\bibitem{PHartshorne}
R.~Hartshorne,
\emph{Algebraic Geometry}, Graduate Texts in Mathematics 52,
Springer, 1977.
\href{https://doi.org/10.1007/978-1-4757-3849-0}{doi:10.1007/978-1-4757-3849-0}.

\bibitem{PEtingofLie}
P.~Etingof,
\emph{Lie Groups and Lie Algebras}, MIT lecture notes, 2024 course version.
\href{https://live.ocw.mit.edu/courses/18-755-lie-groups-and-lie-algebras-ii-spring-2024/mit18_755_s24_lec_full.pdf}{MIT OpenCourseWare 公開原稿}.
本文ではTheorem~11.16, p.~66、Proposition~16.9, p.~88、
Lemma~19.16, p.~104、Corollary~19.18とTheorem~19.19, pp.~105--106、
Definition~21.13, p.~114、および\S32.2を用いる。

\bibitem{PDerdzinskiGal}
A.~Derdzinski and \'{S}.~R.~Gal,
\emph{Curvature spectra of simple Lie groups},
Abhandlungen aus dem Mathematischen Seminar der Universit\"at Hamburg
\textbf{83} (2013), no.~2, 219--230.
\url{https://doi.org/10.1007/s12188-013-0085-z};
\href{https://arxiv.org/abs/1304.2801}{arXiv:1304.2801}.
参照箇所は著者最終原稿の(1.4), p.~2、Theorem~B, p.~3、
Lemma~7.1(a), p.~10。
\url{https://people.math.osu.edu/derdzinski.1/preprints/amsuh14.pdf}.

\bibitem{PBondal93}
A.~I.~Bondal,
\emph{Non-commutative deformations and Poisson brackets on projective spaces},
Max-Planck-Institut f\"ur Mathematik Preprint 93-67 (1993).
\S4、印刷pp.~11--12のEuler--divergence恒等式と無番号Theoremを参照。
\url{https://archive.mpim-bonn.mpg.de/id/eprint/1398/1/preprint_1993_67.pdf}.

\bibitem{PPymElliptic}
B.~Pym,
\emph{Elliptic singularities on log symplectic manifolds and
Feigin--Odesskii Poisson brackets},
Compositio Mathematica \textbf{153} (2017), no.~4, 717--744.
\url{https://doi.org/10.1112/S0010437X16008174};
\href{https://arxiv.org/abs/1507.05668v1}{arXiv:1507.05668v1}.
本文の参照版はarXiv v1：
Definition~4.2, p.~17、Theorem~1.3, p.~4、
Lemma~5.5, pp.~28--29、Proposition~5.7, p.~30。

\bibitem{PBricalliFavalePirola}
D.~Bricalli, F.~F.~Favale and G.~P.~Pirola,
\emph{On the irreducibility of Hessian loci of cubic hypersurfaces},
Forum of Mathematics, Sigma \textbf{13} (2025), e80, 1--36.
\href{https://doi.org/10.1017/fms.2025.36}{doi:10.1017/fms.2025.36};
\href{https://arxiv.org/abs/2406.12024}{arXiv:2406.12024}.
参照箇所は刊行版Theorem~A, p.~2、Theorem~B, p.~3、
Theorem~3.4とその証明, pp.~10--12。

\bibitem{PMayanskiy}
E.~Mayanskiy,
Poisson cohomology of two Fano threefolds,
\emph{Journal of Algebra} \textbf{424} (2015), 21--45.
\href{https://doi.org/10.1016/j.jalgebra.2014.08.049}{doi:10.1016/j.jalgebra.2014.08.049};
\href{https://arxiv.org/abs/1303.5916}{arXiv:1303.5916}.
本文で用いるGrassmann記述はTheorem~1（原稿p.~2）。
同定理および原稿p.~3は分類をLoray--Pereira--Touzetに帰属させている。

\bibitem{BondalGualtieriPym}
M.~Gualtieri and B.~Pym,
\emph{Poisson modules and degeneracy loci},
Proc. Lond. Math. Soc. \textbf{107} (2013), 627--654.
\href{https://doi.org/10.1112/plms/pds090}{doi:10.1112/plms/pds090};
\href{https://arxiv.org/abs/1203.4293}{arXiv:1203.4293}.
Bondal予想の定式化、四次元Fanoの場合の証明。

\bibitem{BondalDruelPereiraPymTouzet}
S.~Druel, J.~V.~Pereira, B.~Pym and F.~Touzet,
\emph{Bondal's conjecture in dimension five},
\href{https://arxiv.org/abs/2606.03424}{arXiv:2606.03424} (2026).
定理番号はarXiv v1による：Theorem~1.2（奇数次元のsubmaximal退化軌跡）、
Corollary~1.3（Fano五次元）。
\url{https://druel.perso.math.cnrs.fr/textes/bondal5d.pdf}.

\bibitem{FuHwangCI}
B. Fu and J.-M. Hwang,
\emph{Isotrivial VMRT-structures of complete intersection type},
Asian Journal of Mathematics \textbf{22} (2018), 333--354.
Lemma~4.1（Br\"uckmannの消滅定理を引用）。
\url{https://arxiv.org/abs/1608.00846v2}.
\url{https://doi.org/10.4310/AJM.2018.v22.n2.a10}.

\bibitem{BorelRemmert}
A. Borel and R. Remmert,
\emph{\"Uber kompakte homogene K\"ahlersche Mannigfaltigkeiten},
Mathematische Annalen \textbf{145} (1962), 429--439, Theorem I.
英訳：\url{https://wglobke.github.io/publications/BorelRemmert1962.pdf}.

\bibitem{LPTPoissonCubic}
F.~Loray, J.~V.~Pereira and F.~Touzet,
\emph{Foliations with trivial canonical bundle on Fano 3-folds},
Math. Nachr. \textbf{286} (2013), 921--940.
\href{https://arxiv.org/abs/1112.3871}{arXiv:1112.3871}.

\bibitem{PStacks}
The Stacks Project Authors,
\emph{The Stacks Project},
Cohen--Macaulay環・Koszul複体・射影空間のコホモロジー・商スタック。
参照タグ：
\href{https://stacks.math.columbia.edu/tag/00ND}{00ND},
\href{https://stacks.math.columbia.edu/tag/02JN}{02JN},
\href{https://stacks.math.columbia.edu/tag/062F}{062F},
\href{https://stacks.math.columbia.edu/tag/01XT}{01XT},
\href{https://stacks.math.columbia.edu/tag/06FI}{06FI}.

\bibitem{PGLefschetz}
A.~Grothendieck,
\emph{Cohomologie locale des faisceaux coh\'erents et th\'eor\`emes
de Lefschetz locaux et globaux (SGA~2)},
Expos\'e~XII, Corollaire~3.7 (Picard groups of sufficiently
high-dimensional smooth projective complete intersections).
\url{https://grothendiecksga.com/read/sga2/en/XII-3.html}.

\bibitem{PMatsumuraMonsky}
H.~Matsumura and P.~Monsky,
\emph{On the automorphisms of hypersurfaces},
J. Math. Kyoto Univ. \textbf{3} (1963/64), 347--361.
\url{https://doi.org/10.1215/kjm/1250524785}.

\bibitem{CubicQuantization}
\emph{滑らかな三次超曲面の代数的量子化},
研究原稿、2026年11月21日。
\href{https://github.com/uedakazushi/cubic-quantization/blob/main/papers/cubic-quantization.tex}{GitHub公開ソース（主ファイル）}.
本稿のPoisson正規形を入力とする代数的中心元・PBW性・同変降下を扱う独立論文。
\end{thebibliography}
